\documentclass{article}

\usepackage{float}
\restylefloat{table}

\usepackage{booktabs}

\title{Team Productivity: POC\\\progname}

\author{\authname}

\date{}

\input{../Comments.text}
\input{../Common.text}

\begin{document}

\maketitle
\sloppy

This document summarizes the contributions of each team member up to the POC
Demo.  The time period of interest is the time between the beginning of the term
and the POC demo.

\noindent \textbf{Project repository:}
\url{https://github.com/shivyam/cas-course-optimization}

\section{Demo Plans}

The main feasibility risk for this project is extracting reliable information
from CAS course outlines. Outlines do not follow a single format, and some
information may be missing or incomplete (for example, ``Midterm Date: TBD'').
Every other feature (the shared calendar, conflict detection, scheduling
suggestions and compliance checking) depends on this information, so if
assessment dates or weights are extracted incorrectly, the rest of the system
will also be unreliable.

\bigskip
\noindent The PoC demonstration shows one course outline going through the
whole extraction pipeline:

\begin{enumerate}
  \item A PDF course outline is uploaded through the web application.
  \item The backend extracts the text from the PDF with PyMuPDF.
  \item A locally hosted open-weight model, run through Ollama, returns JSON in
  a fixed schema: the name, type, date (or TBD) and weight of each assessment,
  the source passage for each assessment, the learning outcomes, and a list of
  missing fields.
  \item The backend checks that the JSON matches the schema and that each
  source passage appears in the extracted text.
  \item The frontend displays the assessments in a table, with the source
  passage beside each value and missing or TBD values clearly flagged.
\end{enumerate}

\noindent We will also report the pipeline's accuracy on a small set of real
outlines scored against our answer key, and show how the pipeline behaves when
an instruction is inserted into an outline (prompt injection).

\section{Team Meeting Attendance}

\begin{table}[H]
\centering
\begin{tabular}{ll}
\toprule
\textbf{Student} & \textbf{Meetings}\\
\midrule
Total & 12\\
Harmanpreet Sagar & 12\\
Shivya Mehta & 12\\
Anindita Bornomala & 12\\
Vanessa Lai & 12\\
\bottomrule
\end{tabular}
\end{table}

The team meets once before and once after each of the six supervisor and stakeholder meetings, to prepare and then to follow up on the action items. All four team members attended every team meeting.

\section{Supervisor/Stakeholder Meeting Attendance}

\noindent \textbf{Supervisor's Name: } Matt Vonk and Connie Carrabs

\begin{table}[H]
\centering
\begin{tabular}{ll}
\toprule
\textbf{Student} & \textbf{Meetings}\\
\midrule
Total & 6\\
Harmanpreet Sagar & 6\\
Shivya Mehta & 6\\
Anindita Bornomala & 6\\
Vanessa Lai & 6\\
\bottomrule
\end{tabular}
\end{table}

The six meetings were four supervisor meetings with Matt Vonk and Connie Carrabs
(October 2 \href{https://github.com/shivyam/cas-course-optimization/issues/25}{\#25}, October 20 \href{https://github.com/shivyam/cas-course-optimization/issues/179}{\#179}, November 3 \href{https://github.com/shivyam/cas-course-optimization/issues/181}{\#181} and
November 17 \href{https://github.com/shivyam/cas-course-optimization/issues/182}{\#182}) and two stakeholder meetings with the CAS Associate
Chairs for Undergraduate Studies (Dr.~Smith on October 6 \href{https://github.com/shivyam/cas-course-optimization/issues/26}{\#26} and
Dr.~Mosser on October 26 \href{https://github.com/shivyam/cas-course-optimization/issues/180}{\#180}). All four members attended every meeting.

\section{Lecture Attendance}

\wss{For each team member how many lectures have they attended over the time
period of interest.  This number should be determined from the lecture issues in
the team's repo. You can find the number of lectures in the time period of
interest by looking at the
\href{https://calendar.google.com/calendar/u/0/embed?src=rnboqiaki1k2la7rpu3bn0um58@group.calendar.google.com&ctz=America/Toronto}
{Google calendar} for the capstone course. The first entry in the
table should be the total number of lectures held for the class during the period of interest.}

\wss{NOTE: There will be approximately 13 lectures between the start of class
and the POC demos}

\begin{table}[H]
\centering
\begin{tabular}{ll}
\toprule
\textbf{Student} & \textbf{Lectures}\\
\midrule
Total & Num\\
Name 1 & Num\\
Name 2 & Num\\
Name 3 & Num\\
Name 4 & Num\\
Name 5 & Num\\
\bottomrule
\end{tabular}
\end{table}

\wss{If needed, an explanation for the lecture attendance can be provided here.}

\section{TA Document Discussion Attendance}

\wss{For each team member how many of the informal document discussion meetings
with the TA were attended over the time period of interest. The first row is the
total number of discussion meetings scheduled by the TA with the team.}

\noindent \textbf{TA's Name: } [fill in this information]

\begin{table}[H]
\centering
\begin{tabular}{ll}
\toprule
\textbf{Student} & \textbf{Meetings}\\
\midrule
Total & Num\\
Name 1 & Num\\
Name 2 & Num\\
Name 3 & Num\\
Name 4 & Num\\
Name 5 & Num\\
\bottomrule
\end{tabular}
\end{table}

\wss{If needed, an explanation for the attendance can be provided here.}

\section{Commits, Lines Added and Lines Deleted, Pull Requests}

\wss{For each team member how many commits to the main branch have been made
over the time period of interest.  The total is the total number of commits for
the entire team since the beginning of the term.  The percentage is the
percentage of the total commits made by each team member.}

\wss{There are also columns for lines added and lines deleted to give a sense of
the magnitude of each team member's contributions.  Note that these numbers can
be somewhat misleading, as a team member who does refactoring may have high
numbers in both columns, while a team member who implements a large new feature
may have a high number of lines added, but few lines deleted.  As for all commit
information, these numbers will be interpreted with care.}

\wss{The pull requests column counts the number of pull requests created by each
team member.  This is a measure of how many contributions were made in branches,
which may not be reflected in the commit counts for the main branch.}

\begin{table}[H]
\centering
\begin{tabular}{lllllll}
\toprule
\textbf{Student} & \textbf{Commits} & \textbf{Percent} & \textbf{Lines Added} &
\textbf{Lines Deleted} & \textbf{Pull Requests}\\
\midrule
Total & Num & 100\% & total added & total deleted & total PRs\\
Name 1 & Num & \% & num & num & num\\
Name 2 & Num & \% & num & num & num\\
Name 3 & Num & \% & num & num & num\\
Name 4 & Num & \% & num & num & num\\
Name 5 & Num & \% & num & num & num\\
\bottomrule
\end{tabular}
\end{table}

\wss{If needed, an explanation for the counts can be provided here.  For
instance, if a team member has more commits to unmerged branches, these numbers
can be provided here.  If multiple people contribute to a commit, git allows for
multi-author commits.}

\section{Issue Tracker}

\wss{For each team member how many issues have they authored (including open and
closed issues (O+C)) and how many have they been assigned (only counting closed
issues (C only)) over the time period of interest.}

\begin{table}[H]
\centering
\begin{tabular}{lll}
\toprule
\textbf{Student} & \textbf{Authored (O+C)} & \textbf{Assigned (C only)}\\
\midrule
Name 1 & Num & Num \\
Name 2 & Num & Num \\
Name 3 & Num & Num \\
Name 4 & Num & Num \\
Name 5 & Num & Num \\
\bottomrule
\end{tabular}
\end{table}

\wss{If needed, an explanation for the counts can be provided here.}

\section{CICD via GitHub Actions}

\wss{Say how CICD technology is used and will be used in your project.}

\wss{Provide links to your CICD yaml files}

\section{Custom Documents (CDs)}

\wss{What is the plan (as documented in TeamComposition.csv) for the team's
CDs?  Should the CDs be modified now that the team knows more about the
project?}

\section{Team Charter Trigger Items}

\wss{Provide a summary of the quantified triggers identified in the team's
charter.}

\wss{Provide a list of any violations of the triggers.  If the team wishes, the
violations can be summarized on aggregate, instead of naming specific team
members.}

\wss{Provide a plan to address the violations.  This could include revising the
triggers, if they are found to be too weak, strong or ambiguous.}

\section{Additional Productivity Metrics}

\wss{If your team has additional metrics of productivity, please feel free to
add them to this report.  If not, please explicitly state that there are no additional
metrics.}

\wss{Additional metrics can include things like code reviews done, pull requests
created, count of joining meetings late, count of number of times contributions
had to be corrected, number of internal deadlines missed, test cases written, etc.}

\wss{We are looking for data on these metrics, not just a list of additional
metrics the team is planning on tracking.  However, if all you have is a plan,
please share it here.}

\end{document}